% GENERATED FILE. Edit canonical lessons, then run npm run generate:books.
% canonical-source-hash: fnv1a64:376ad63182c7ecb3
% canonical-lessons: KA-C160-baggu, KA-C160-horu, KA-C160-vyayama-madu, KA-C160-sinu, KA-C160-nadugu

\chapter{To bend down, To carry, To exercise, To sneeze, To shiver}
\label{ch:ka-to-bend-down-to-carry-to-exercise-to-sneeze-to-shiver}
\hlchaptermodality{160}

\begin{chapteropening}
\textbf{By the end of this chapter:} \emph{I can say to bend down, to carry, to exercise, to sneeze and to shiver.}

Complete the last lesson of chapter 160: I can say to bend down, to carry, to exercise, to sneeze and to shiver.
\end{chapteropening}

\section[\texorpdfstring{\kn{ಬಗ್ಗು} (baggu)}{ಬಗ್ಗು (baggu)}]{\kn{ಬಗ್ಗು} (baggu) — to bend down}

\label{lesson:KA-C160-baggu}



\begin{quote}
Before the new one: say the Kannada for to taste, then the Kannada for to chew.
\end{quote}

\subsection*{You'll want to know: \kn{ಬಗ್ಗು}}
\textbf{\kn{ಬಗ್ಗು}} — \emph{baggu} — \textquotedblleft{}to bend down\textquotedblright{}.

\begin{grammarlens}[title={What you've built}]
One more word for this chapter.
\end{grammarlens}

\subsection*{Your turn}
\begin{itemize}
\raggedright
  \item \emph{Say it:} \emph{baggu}
  \item \emph{Say it:} \emph{baggu}, once more
  \item \emph{Say it:} say \emph{baggu}
  \item \emph{Recall:} say the Kannada for to taste, then the Kannada for to chew, then say \emph{baggu} again
\end{itemize}

\begin{tcolorbox}[breakable,colback=teal!4,colframe=teal!35!black,title={Before you move on}]
What does \kn{ಬಗ್ಗು} mean? (\textquotedblleft{}To bend down\textquotedblright{}.) Say it once more.
\end{tcolorbox}
\section[\texorpdfstring{\kn{ಹೊರು} (horu)}{ಹೊರು (horu)}]{\kn{ಹೊರು} (horu) — to carry (a load)}

\label{lesson:KA-C160-horu}



\begin{quote}
Before the new one: say the Kannada for to chew, then the Kannada for to bend down.
\end{quote}

\subsection*{You'll want to know: \kn{ಹೊರು}}
\textbf{\kn{ಹೊರು}} — \emph{horu} — \textquotedblleft{}to carry (a load)\textquotedblright{}.

\begin{grammarlens}[title={What you've built}]
One more word for this chapter.
\end{grammarlens}

\subsection*{Your turn}
\begin{itemize}
\raggedright
  \item \emph{Say it:} \emph{horu}
  \item \emph{Say it:} \emph{horu}, once more
  \item \emph{Say it:} say \emph{horu}
  \item \emph{Recall:} say the Kannada for to chew, then the Kannada for to bend down, then say \emph{horu} again
\end{itemize}

\begin{tcolorbox}[breakable,colback=teal!4,colframe=teal!35!black,title={Before you move on}]
What does \kn{ಹೊರು} mean? (\textquotedblleft{}To carry (a load)\textquotedblright{}.) Say it once more.
\end{tcolorbox}
\section[\texorpdfstring{\kn{ವ್ಯಾಯಾಮ} \kn{ಮಾಡು} (vyāyāma māḍu)}{ವ್ಯಾಯಾಮ ಮಾಡು (vyāyāma māḍu)}]{\kn{ವ್ಯಾಯಾಮ} \kn{ಮಾಡು} (vyāyāma māḍu) — to exercise}

\label{lesson:KA-C160-vyayama-madu}



\begin{quote}
Before the new one: say the Kannada for to bend down, then the Kannada for to carry.
\end{quote}

\subsection*{You'll want to know: \kn{ವ್ಯಾಯಾಮ} \kn{ಮಾಡು}}
\textbf{\kn{ವ್ಯಾಯಾಮ} \kn{ಮಾಡು}} — \emph{vyāyāma māḍu} — \textquotedblleft{}to exercise\textquotedblright{}.

\begin{grammarlens}[title={What you've built}]
One more word for this chapter.
\end{grammarlens}

\subsection*{Your turn}
\begin{itemize}
\raggedright
  \item \emph{Say it:} \emph{vyāyāma māḍu}
  \item \emph{Say it:} \emph{vyāyāma māḍu}, once more
  \item \emph{Say it:} say \emph{vyāyāma māḍu}
  \item \emph{Recall:} say the Kannada for to bend down, then the Kannada for to carry, then say \emph{vyāyāma māḍu} again
\end{itemize}

\begin{tcolorbox}[breakable,colback=teal!4,colframe=teal!35!black,title={Before you move on}]
What does \kn{ವ್ಯಾಯಾಮ} \kn{ಮಾಡು} mean? (\textquotedblleft{}To exercise\textquotedblright{}.) Say it once more.
\end{tcolorbox}
\section[\texorpdfstring{\kn{ಸೀನು} (sīnu)}{ಸೀನು (sīnu)}]{\kn{ಸೀನು} (sīnu) — to sneeze}

\label{lesson:KA-C160-sinu}



\begin{quote}
Before the new one: say the Kannada for to carry, then the Kannada for to exercise.
\end{quote}

\subsection*{You'll want to know: \kn{ಸೀನು}}
\textbf{\kn{ಸೀನು}} — \emph{sīnu} — \textquotedblleft{}to sneeze\textquotedblright{}.

\begin{grammarlens}[title={What you've built}]
One more word for this chapter.
\end{grammarlens}

\subsection*{Your turn}
\begin{itemize}
\raggedright
  \item \emph{Say it:} \emph{sīnu}
  \item \emph{Say it:} \emph{sīnu}, once more
  \item \emph{Say it:} say \emph{sīnu}
  \item \emph{Recall:} say the Kannada for to carry, then the Kannada for to exercise, then say \emph{sīnu} again
\end{itemize}

\begin{tcolorbox}[breakable,colback=teal!4,colframe=teal!35!black,title={Before you move on}]
What does \kn{ಸೀನು} mean? (\textquotedblleft{}To sneeze\textquotedblright{}.) Say it once more.
\end{tcolorbox}
\section[\texorpdfstring{\kn{ನಡುಗು} (naḍugu)}{ನಡುಗು (naḍugu)}]{\kn{ನಡುಗು} (naḍugu) — to shiver, to tremble}

\label{lesson:KA-C160-nadugu}



\begin{quote}
Before the new one: say the Kannada for to exercise, then the Kannada for to sneeze.
\end{quote}

\subsection*{You'll want to know: \kn{ನಡುಗು}}
\textbf{\kn{ನಡುಗು}} — \emph{naḍugu} — \textquotedblleft{}to shiver, to tremble\textquotedblright{}.

\begin{grammarlens}[title={What you've built}]
One more word for this chapter.
\end{grammarlens}

\subsection*{Your turn}
\begin{itemize}
\raggedright
  \item \emph{Say it:} \emph{naḍugu}
  \item \emph{Say it:} \emph{naḍugu}, once more
  \item \emph{Say it:} say \emph{naḍugu}
  \item \emph{Recall:} say the Kannada for to exercise, then the Kannada for to sneeze, then say \emph{naḍugu} again
\end{itemize}

\begin{tcolorbox}[breakable,colback=teal!4,colframe=teal!35!black,title={Before you move on}]
What does \kn{ನಡುಗು} mean? (\textquotedblleft{}To shiver, to tremble\textquotedblright{}.) Say it once more.
\end{tcolorbox}
